\documentclass[11pt]{article}

\usepackage[a4paper,margin=29mm]{geometry}
\usepackage[T1]{fontenc}
\usepackage[utf8]{inputenc}
\usepackage{lmodern}
\usepackage{microtype}
\usepackage{amsmath,amssymb,amsthm,mathtools}
\usepackage{booktabs}
\usepackage{array}
\usepackage{graphicx}
\usepackage{xcolor}
\usepackage{enumitem}
\usepackage{listings}
\usepackage{hyperref}
\usepackage[nameinlink,capitalise,noabbrev]{cleveref}

\definecolor{navy}{RGB}{26,55,91}
\definecolor{softblue}{RGB}{235,242,250}
\definecolor{softred}{RGB}{252,239,239}
\hypersetup{
  colorlinks=true,
  linkcolor=navy,
  citecolor=navy,
  urlcolor=navy,
  pdftitle={A One-Point Resolvent--Hausdorff Criterion for the Riemann Xi Function},
  pdfauthor={Lluís Eriksson}
}

\lstdefinestyle{lean}{
  basicstyle=\ttfamily\small,
  frame=single,
  breaklines=true,
  columns=fullflexible,
  backgroundcolor=\color{softblue},
  showstringspaces=false
}

\newtheorem{theorem}{Theorem}[section]
\newtheorem{proposition}[theorem]{Proposition}
\newtheorem{lemma}[theorem]{Lemma}
\newtheorem{corollary}[theorem]{Corollary}
\theoremstyle{definition}
\newtheorem{definition}[theorem]{Definition}
\newtheorem{remark}[theorem]{Remark}
\newtheorem{researchtarget}[theorem]{Research target}

\newcommand{\C}{\mathbb C}
\newcommand{\R}{\mathbb R}
\newcommand{\N}{\mathbb N}
\newcommand{\XiR}{\mathcal S_{\Xi}}
\newcommand{\RH}{\mathrm{RH}}
\newcommand{\slit}{\C\setminus(-\infty,0]}
\newcommand{\dd}{\,\mathrm d}
\newcommand{\Tr}{\operatorname{Tr}}
\newcommand{\arsinh}{\operatorname{arsinh}}

\title{\textbf{A One-Point Resolvent--Hausdorff Criterion\\for the Riemann $\Xi$ Function}}
\author{Lluís Eriksson}
\date{Research artifact v0.2.0 --- 22 June 2026}

\begin{document}
\maketitle

\begin{center}
\fcolorbox{navy}{softred}{\parbox{0.91\textwidth}{
\textbf{Status and scope.} This manuscript does not claim a proof of the Riemann
hypothesis. It proves an abstract reduction, develops finite certificate hierarchies,
and supplies a finite Lean 4 proof layer with mandatory CI verification. The concrete spectral convergence
that would imply RH remains open. Novelty relative to existing Hausdorff-moment criteria
also requires specialist review.
}}
\end{center}

\begin{abstract}
For $x>1/4$ define
\[
 \XiR(x)=\frac{1}{2\sqrt{x}}\frac{\xi'}{\xi}
 \left(\frac12+\sqrt{x}\right),
 \qquad
 \Xi(z)=\xi\left(\frac12+iz\right).
\]
We isolate a slit-plane extension criterion: if $\XiR$ agrees on one real interval
with a holomorphic function on $\C\setminus(-\infty,0]$, then every zero of $\Xi$
is real. Fixing a single $x_0>1/4$, we encode the derivative jet by
\[
 b_n(x_0)=x_0^n\frac{(-1)^n}{n!}\XiR^{(n)}(x_0).
\]
Using the Hausdorff moment theorem, we show that RH is equivalent to the sequence
$(b_n(x_0))_{n\ge0}$ being a moment sequence of a finite positive measure on $[0,1]$,
or equivalently to all signed forward differences being nonnegative. Under RH the
measure is obtained by the resolvent compactification
$\gamma^2\mapsto x_0/(\gamma^2+x_0)$ of the ordinates of the zeros. Conversely, a
Hausdorff representing measure reconstructs a slit-plane extension. We derive finite
difference and positive semidefinite certificate hierarchies, an explicit prime-side
target with a uniform elementary truncation bound, and a normal-family criterion for
positive squared-resolvent approximants. A Lean 4 companion encodes proofs of the finite atomic
identities, compactification inequalities, sum-of-squares certificates, and error-budget
glue. The decisive open task is a non-circular convergence estimate for a concrete
prime-built self-adjoint family.
\end{abstract}

\tableofcontents

\section{Introduction}

The Riemann hypothesis is the assertion that every non-trivial zero of the zeta
function has real part $1/2$ \cite{Riemann1859,Titchmarsh1986,ClayRH}. A classical
spectral ambition is to realize the corresponding ordinates as the spectrum of a
self-adjoint operator. Recent work has produced prime-built self-adjoint approximants
with striking numerical behaviour, while identifying rigorous convergence as the
central unresolved step \cite{ConnesConsaniMoscovici2025,ConnesVanSuijlekom2025}.

This paper advocates a positive scalar observable rather than direct convergence of
regularized determinants:
\[
 S_D(x)=\frac12\Tr(D^2+xI)^{-1}.
\]
It suppresses high spectral values, is stable under perturbation, and belongs to the
Stieltjes class whenever the squared resolvent is trace class. The transformation
\[
 \lambda\longmapsto \frac{x_0}{\lambda+x_0}
\]
compactifies a nonnegative spectrum to $[0,1]$. At a single fixed point $x_0$, the
entire derivative jet of the target therefore becomes a candidate Hausdorff moment
sequence.

Moment formulations of RH already exist; in particular Zhang develops Hausdorff-moment
criteria for classes of entire functions and for RH \cite{Zhang2023}. Accordingly, the
present manuscript does not assert that ``moments imply RH'' is new. Its proposed
contribution is the integrated package: a squared-resolvent target, one-point derivative
jet, slit-plane reconstruction, prime-side truncation, finite semidefinite certificates,
and formal verification. Whether the precise one-point formulation is novel is a
literature question, not a theorem.

\paragraph{Contributions.}
The paper establishes the following chain.
\begin{enumerate}[label=(\roman*)]
  \item A slit-plane extension of $\XiR$ excludes non-real zeros of $\Xi$.
  \item For every fixed $x_0>1/4$, RH is equivalent to a Hausdorff moment condition on
  the derivatives of $\XiR$ at $x_0$.
  \item The moment condition yields scalar finite-difference and finite PSD tests.
  \item The arithmetic target has an absolutely convergent prime expansion and an
  explicit elementary tail bound.
  \item Positive Stieltjes approximants with one-point control need only converge on one
  real interval to imply RH.
  \item A Lean 4 artifact encodes the finite algebraic substrate without project axioms and makes compilation an explicit CI gate.
\end{enumerate}

\section{The completed zeta function and the resolvent target}

Define
\begin{equation}\label{eq:xi-def}
 \xi(s)=\frac12s(s-1)\pi^{-s/2}\Gamma\left(\frac{s}{2}\right)\zeta(s),
 \qquad
 \Xi(z)=\xi\left(\frac12+iz\right).
\end{equation}
The function $\xi$ is entire and satisfies $\xi(s)=\xi(1-s)$; hence $\Xi$ is entire,
even, and real on the real axis. The non-trivial zero $\rho=\beta+i\gamma$ of $\zeta$
corresponds to the zero
\[
 z=\gamma-i\left(\beta-\frac12\right)
\]
of $\Xi$. Thus
\begin{equation}\label{eq:rh-real-zeros}
 \RH\quad\Longleftrightarrow\quad\text{all zeros of $\Xi$ are real}.
\end{equation}

For $x>1/4$, set
\begin{equation}\label{eq:target}
 \XiR(x)=\frac{1}{2\sqrt{x}}\frac{\xi'}{\xi}
 \left(\frac12+\sqrt{x}\right).
\end{equation}
The restriction ensures $\sigma=1/2+\sqrt{x}>1$, where the Euler product is absolutely
convergent and $\zeta(\sigma)\ne0$.

\begin{proposition}[Spectral representation under RH]\label{prop:rh-stieltjes}
Assume RH and enumerate the positive zeros of $\Xi$, with multiplicity, by
$0<\gamma_1\le\gamma_2\le\cdots$. Then
\begin{equation}\label{eq:stieltjes-zeros}
 \XiR(x)=\sum_{j\ge1}\frac{1}{\gamma_j^2+x},\qquad x>0.
\end{equation}
In particular $\XiR$ is the Stieltjes transform of the positive discrete measure
$\sum_j\delta_{\gamma_j^2}$.
\end{proposition}

\begin{proof}
Because $\Xi$ is even of order one and $\sum_j\gamma_j^{-2}<\infty$, its symmetric
canonical product has the form
\[
 \Xi(z)=\Xi(0)\prod_{j\ge1}\left(1-\frac{z^2}{\gamma_j^2}\right).
\]
Logarithmic differentiation gives
\[
 -\frac{\Xi'(z)}{\Xi(z)}=2z\sum_{j\ge1}\frac{1}{\gamma_j^2-z^2}.
\]
At $z=iy$, the functional equation implies
\[
 -\frac{\Xi'(iy)}{\Xi(iy)}
 =i\frac{\xi'}{\xi}\left(\frac12+y\right).
\]
Dividing by $2iy$ and putting $x=y^2$ yields \eqref{eq:stieltjes-zeros}. Standard zero
counting gives convergence locally uniformly away from the negative real cut.
\end{proof}

\section{A slit-plane extension criterion}

Let
\[
 \Omega=\slit.
\]
The quadratic map used below is the key geometric device.

\begin{lemma}\label{lem:quadratic-map}
If $\operatorname{Im}z>0$, then $-z^2\in\Omega$.
\end{lemma}

\begin{proof}
Write $z=a+ib$ with $b>0$. If $a=0$, then $-z^2=b^2>0$. If $a\ne0$, then
$\operatorname{Im}(-z^2)=-2ab\ne0$, so $-z^2$ is not on the non-positive real axis.
\end{proof}

\begin{theorem}[Slit-plane extension criterion]\label{thm:slit-extension}
Let $I\subset(1/4,\infty)$ be a nonempty open interval. Suppose that a holomorphic
function $S:\Omega\to\C$ satisfies
\[
 S(x)=\XiR(x),\qquad x\in I.
\]
Then RH holds.
\end{theorem}

\begin{proof}
For $z$ in the upper half-plane define
\[
 M(z)=2zS(-z^2).
\]
By \cref{lem:quadratic-map}, $M$ is holomorphic. If $z=iy$ and $y^2\in I$, then
\[
 M(iy)=2iy\XiR(y^2)
 =i\frac{\xi'}{\xi}\left(\frac12+y\right)
 =-\frac{\Xi'(iy)}{\Xi(iy)}.
\]
The meromorphic identity theorem therefore gives
\[
 M(z)=-\frac{\Xi'(z)}{\Xi(z)}
\]
throughout the upper half-plane away from the zeros of $\Xi$.

Suppose $z_0$ were an upper-half-plane zero of multiplicity $m>0$. Locally
$\Xi(z)=(z-z_0)^m g(z)$ with $g(z_0)\ne0$, and hence
\[
 -\frac{\Xi'(z)}{\Xi(z)}=-\frac{m}{z-z_0}-\frac{g'(z)}{g(z)}.
\]
This has a non-removable pole at $z_0$, contradicting the holomorphy of $M$. Thus no
zero lies in the upper half-plane. Conjugation symmetry excludes the lower half-plane,
so all zeros are real. By \eqref{eq:rh-real-zeros}, RH follows.
\end{proof}

\begin{remark}
The theorem does not reconstruct every zero, require determinant convergence, or identify
an inverse spectral problem. Agreement on one interval is enough because analyticity
propagates the logarithmic derivative.
\end{remark}

\section{The one-point Hausdorff criterion}

Fix once and for all $x_0>1/4$. Define the normalized derivative jet
\begin{equation}\label{eq:bn-def}
 b_n(x_0)=x_0^n\frac{(-1)^n}{n!}\XiR^{(n)}(x_0),\qquad n\ge0.
\end{equation}
For a sequence $b=(b_n)$ define the signed forward differences recursively by
\[
 D^0b_n=b_n,\qquad D^{k+1}b_n=D^kb_n-D^kb_{n+1}.
\]
Thus $D^k=(-1)^k\Delta^k$ under the usual convention $\Delta b_n=b_{n+1}-b_n$.

We use the classical Hausdorff moment theorem \cite{Hausdorff1921,ShohatTamarkin1943,Akhiezer1965}:
\begin{equation}\label{eq:hausdorff-theorem}
 b_n=\int_0^1 v^n\dd\mu(v)\text{ for a finite positive measure $\mu$}
 \quad\Longleftrightarrow\quad
 D^kb_n\ge0\text{ for all $k,n\ge0$}.
\end{equation}

\begin{theorem}[One-point resolvent--Hausdorff criterion]\label{thm:one-point}
For every fixed $x_0>1/4$, the following are equivalent.
\begin{enumerate}[label=(\alph*)]
  \item RH holds.
  \item The sequence $\bigl(b_n(x_0)\bigr)_{n\ge0}$ is a Hausdorff moment sequence.
  \item For all $k,n\ge0$,
  \[
    D^k b_n(x_0)\ge0.
  \]
\end{enumerate}
\end{theorem}

\begin{proof}
Assume RH. By \cref{prop:rh-stieltjes}, differentiation term by term gives
\[
 \frac{(-1)^n}{n!}\XiR^{(n)}(x_0)
 =\sum_{j\ge1}\frac{1}{(\gamma_j^2+x_0)^{n+1}}.
\]
Put
\[
 u_j=\frac{1}{\gamma_j^2+x_0},
 \qquad
 v_j=\frac{x_0}{\gamma_j^2+x_0}\in(0,1).
\]
Then
\[
 b_n(x_0)=\sum_{j\ge1}u_jv_j^n
 =\int_0^1 v^n\dd\mu_{x_0}(v),
 \qquad
 \mu_{x_0}=\sum_{j\ge1}u_j\delta_{v_j}.
\]
The measure is finite because its total mass is $\XiR(x_0)$. This proves (a)$\Rightarrow$(b),
and (b)$\Leftrightarrow$(c) is \eqref{eq:hausdorff-theorem}.

Now assume (b), with representing measure $\mu$. For $w\in\Omega$ define
\begin{equation}\label{eq:hausdorff-extension}
 F(w)=\int_0^1
 \frac{\dd\mu(v)}{1+\frac{w-x_0}{x_0}v}.
\end{equation}
For $v>0$, the only zero of the denominator occurs at
$w=x_0(1-v^{-1})\in(-\infty,0]$; for $v=0$ the denominator is one. Standard dominated
holomorphic-parameter integration shows that $F$ is holomorphic on $\Omega$.

For $h$ sufficiently small,
\[
 F(x_0+h)
 =\int_0^1\sum_{n\ge0}\left(-\frac{h}{x_0}v\right)^n\dd\mu(v)
 =\sum_{n\ge0}(-1)^n\frac{h^n}{x_0^n}b_n(x_0).
\]
By \eqref{eq:bn-def}, this is precisely the Taylor series of $\XiR(x_0+h)$.
Consequently $F$ agrees with $\XiR$ on a nonempty interval around $x_0$. Applying
\cref{thm:slit-extension} proves RH.
\end{proof}

\begin{corollary}[Explicit finite differences under RH]\label{cor:diff-spectral}
Under RH,
\[
 D^k b_n(x_0)
 =\sum_{j\ge1}
 \frac{1}{\gamma_j^2+x_0}
 \left(\frac{x_0}{\gamma_j^2+x_0}\right)^n
 \left(1-\frac{x_0}{\gamma_j^2+x_0}\right)^k
 \ge0.
\]
\end{corollary}

\section{Finite certificate hierarchies}

The infinite family in \cref{thm:one-point} can be truncated in several falsifiable ways.
A single certified negative value of $D^kb_n$ would refute RH. Passing a finite number of
tests is evidence about the truncation only, never a proof of RH.

\subsection{Finite differences}

For integers $K,N\ge0$, define the scalar test region
\[
 \mathcal D_{K,N}=\{D^kb_n:0\le k\le K,\ 0\le n\le N\}.
\]
Each quantity can in principle be evaluated from derivatives of the prime-side target, with
interval arithmetic used to control truncation and differentiation errors.

\subsection{Hankel and localizing forms}

Given a moment sequence $(b_n)$, define
\[
 H_N=(b_{i+j})_{0\le i,j<N},
 \qquad
 L_N=(b_{i+j}-b_{i+j+1})_{0\le i,j<N}.
\]
If $b_n=\int_0^1v^n\dd\mu(v)$, then for every $c\in\R^N$,
\begin{align*}
 c^\mathsf{T}H_Nc
 &=\int_0^1\left|\sum_{i=0}^{N-1}c_iv^i\right|^2\dd\mu(v)\ge0,\\
 c^\mathsf{T}L_Nc
 &=\int_0^1(1-v)\left|\sum_{i=0}^{N-1}c_iv^i\right|^2\dd\mu(v)\ge0.
\end{align*}
Thus
\[
 H_N\succeq0,\qquad L_N\succeq0.
\]
The Lean source encodes proofs of these facts for finite atomic measures as explicit sum-of-squares
certificates, avoiding reliance on numerical eigenvalue computations.

\subsection{Resolvent compactification}

Let $\lambda_i\ge0$ be a finite squared spectrum and $x_0>0$. Define
\[
 w_i=\frac{1}{\lambda_i+x_0},
 \qquad
 p_i=\frac{x_0}{\lambda_i+x_0}\in[0,1],
 \qquad
 b_n^{\mathrm{fin}}=\sum_iw_ip_i^n.
\]
Then
\begin{equation}\label{eq:finite-difference-identity}
 D^kb_n^{\mathrm{fin}}
 =\sum_iw_ip_i^n(1-p_i)^k\ge0.
\end{equation}
This identity is the central kernel-checked theorem in the current Lean layer.

\section{The arithmetic target and an explicit prime tail}

For real $\sigma>1$, logarithmic differentiation of \eqref{eq:xi-def} gives
\begin{equation}\label{eq:xi-log-derivative}
 \frac{\xi'}{\xi}(\sigma)
 =\frac1\sigma+\frac1{\sigma-1}-\frac12\log\pi
 +\frac12\psi\left(\frac\sigma2\right)
 -\sum_{m=2}^{\infty}\frac{\Lambda(m)}{m^\sigma},
\end{equation}
where $\psi=\Gamma'/\Gamma$ and $\Lambda$ is the von Mangoldt function. With
$y=\sqrt{x}$ and $\sigma=1/2+y$,
\begin{equation}\label{eq:prime-target}
 \XiR(x)=\frac{1}{2y}\left[
 \frac1\sigma+\frac1{\sigma-1}-\frac12\log\pi
 +\frac12\psi\left(\frac\sigma2\right)
 -\sum_{m=2}^{\infty}\frac{\Lambda(m)}{m^\sigma}
 \right].
\end{equation}
For a cutoff $X\ge3$, let $P_X(x)$ be the same expression with $m\le X$.

\begin{proposition}[Elementary prime-tail bound]\label{prop:prime-tail}
Let $\delta>0$ and $x\ge(1/2+\delta)^2$. For $X\ge3$,
\begin{equation}\label{eq:tail-bound}
 |\XiR(x)-P_X(x)|
 \le
 \frac{X^{-\delta}}{1+2\delta}
 \left(\frac{\log X}{\delta}+\frac{1}{\delta^2}\right).
\end{equation}
\end{proposition}

\begin{proof}
Here $y\ge1/2+\delta$ and $\sigma\ge1+\delta$. Since
$0\le\Lambda(m)\le\log m$,
\[
 |\XiR(x)-P_X(x)|
 \le\frac{1}{2y}\sum_{m>X}\frac{\log m}{m^{1+\delta}}.
\]
For $X\ge3$ the summand is eventually decreasing, and the standard integral majorant gives
\[
 \sum_{m>X}\frac{\log m}{m^{1+\delta}}
 \le\int_X^\infty\frac{\log t}{t^{1+\delta}}\dd t
 =X^{-\delta}\left(\frac{\log X}{\delta}+\frac1{\delta^2}\right).
\]
Finally $(2y)^{-1}\le(1+2\delta)^{-1}$.
\end{proof}

\section{Positive spectral approximants}

Let $\nu_j$ be positive measures on $[0,\infty)$ satisfying
\[
 \int_0^\infty\frac{\dd\nu_j(t)}{t+x_*}<\infty
\]
for some $x_*>0$, and define
\[
 S_j(w)=\int_0^\infty\frac{\dd\nu_j(t)}{t+w},\qquad w\in\Omega.
\]
For a paired discrete spectrum $\{\pm\gamma_{j,k}\}$ this takes the form
\[
 S_j(x)=\sum_k\frac{1}{\gamma_{j,k}^2+x}
 =\frac12\Tr(D_j^2+xI)^{-1}.
\]

\begin{lemma}[One-point normal-family control]\label{lem:one-point-control}
If $\sup_jS_j(x_*)<\infty$, then $(S_j)$ is locally bounded on $\Omega$.
\end{lemma}

\begin{proof}
For every compact $K\subset\Omega$ there exists $C_K<\infty$ such that
\[
 \frac1{|t+w|}\le\frac{C_K}{t+x_*},
 \qquad t\ge0,\ w\in K.
\]
Integration yields $|S_j(w)|\le C_KS_j(x_*)$, uniformly in $j$.
\end{proof}

\begin{theorem}[Abstract prime--resolvent criterion]\label{thm:abstract-resolvent}
Let $I\subset(1/4,\infty)$ be a nonempty open interval. Suppose the $S_j$ above satisfy
\[
 \sup_jS_j(x_*)<\infty
 \qquad\text{and}\qquad
 S_j(x)\longrightarrow\XiR(x)\quad(x\in I).
\]
Then RH holds.
\end{theorem}

\begin{proof}
By \cref{lem:one-point-control} and Montel's theorem, a subsequence converges uniformly on
compact subsets of $\Omega$ to a holomorphic function $S$. Pointwise convergence on $I$
forces $S=\XiR$ there. The conclusion follows from \cref{thm:slit-extension}.
\end{proof}

\begin{corollary}[Prime--resolvent error formulation]\label{cor:prime-resolvent-error}
Fix $\delta>0$ and a compact interval
$I_{\delta,B}\subset((1/2+\delta)^2,\infty)$. Let $X_j\to\infty$. If
\[
 \sup_{x\in I_{\delta,B}}|S_j(x)-P_{X_j}(x)|\to0
\]
and the one-point bound in \cref{thm:abstract-resolvent} holds, then RH holds.
\end{corollary}

\begin{proof}
Combine the assumed comparison with the uniform bound \eqref{eq:tail-bound}, then apply
\cref{thm:abstract-resolvent}.
\end{proof}

\section{A three-part error budget}

For a concrete operator construction, the useful decomposition is
\begin{align*}
 |S_{\mathrm{spec}}-\XiR|
 \le{}&|S_{\mathrm{spec}}-S_{\mathrm{model}}|
 +|S_{\mathrm{model}}-P_X|
 +|P_X-\XiR|.
\end{align*}
The three terms represent:
\begin{enumerate}[label=(\arabic*)]
  \item spectral/Galerkin or state-alignment error;
  \item model normalization and transform-convention error;
  \item the explicit arithmetic tail in \cref{prop:prime-tail}.
\end{enumerate}
Only the third term is closed here in a source-independent manner. For the prime-built
operators of \cite{ConnesConsaniMoscovici2025}, the first two require a precise audit of
domains, Fourier/Mellin conventions, spectral gaps, and trace-class estimates.

The finite-dimensional resolvent identity
\[
 (A^2+xI)^{-1}-(B^2+xI)^{-1}
 =(A^2+xI)^{-1}(B^2-A^2)(B^2+xI)^{-1}
\]
suggests the trace estimate
\[
 \left|\Tr\left((A^2+xI)^{-1}-(B^2+xI)^{-1}\right)\right|
 \le x^{-2}\|A^2-B^2\|_1,
\]
whenever the indicated trace-class hypotheses hold. Establishing a useful version for the
concrete unbounded models is a substantive analytic task, not formal bookkeeping.

\section{Numerical illustration and exact toy certificates}

The accompanying scripts compare \eqref{eq:target}, the prime truncation
\eqref{eq:prime-target}, and a finite sum over the first critical-line zeros. The default
run uses $40$ zeros, prime-power cutoff $20{,}000$, and $70$ points on $[1,25]$.
These choices are intentionally modest: the goal is to test conventions and the artifact
pipeline, not to provide new evidence for RH.

\begin{figure}[ht]
 \centering
 \includegraphics[width=0.88\textwidth]{figures/resolvent_comparison.pdf}
 \caption{Numerical comparison of the completed-zeta target, a truncated von Mangoldt
 expansion, and a partial spectral sum. The visible spectral offset is the omitted positive
 tail, not a contradiction.}
 \label{fig:comparison}
\end{figure}

\begin{figure}[ht]
 \centering
 \includegraphics[width=0.88\textwidth]{figures/resolvent_errors.pdf}
 \caption{Illustrative absolute errors for the same finite calculation. Neither curve is a
 certified infinite-dimensional convergence result.}
 \label{fig:errors}
\end{figure}

The exact certificate script uses
\[
 x_0=1,\qquad \lambda\in\{1,4,9\},
\]
so that
\[
 p\in\left\{\frac12,\frac15,\frac1{10}\right\},
 \qquad w=p.
\]
It exports rational moments, signed differences, Hankel and localizing matrices, and exact
principal minors. Because the support has three atoms, the Hankel rank stabilizes at three;
zero higher principal minors are therefore expected.

\section{Lean 4 formalization}

The artifact is written for Lean 4 and Mathlib \cite{deMoura2015,Mathlib2020}, pinned to
Lean/Mathlib release v4.31.0. The current source-complete layer is deliberately finite. The construction container used for this archive did not contain Lean, so compilation remains an explicit CI release check rather than a locally claimed success.

\begin{table}[ht]
\centering
\small
\begin{tabular}{p{0.31\textwidth}p{0.58\textwidth}}
\toprule
Module & Certified content \\
\midrule
\texttt{Basic} & Atomic moments, signed forward differences, complete-monotonicity predicate. \\
\texttt{HausdorffFinite} & Exact identity \eqref{eq:finite-difference-identity} for finite atoms. \\
\texttt{ResolventCompactification} & Positivity and $[0,1]$ support of compactified spectral atoms. \\
\texttt{FiniteCertificates} & Hankel and localizing sum-of-squares nonnegativity. \\
\texttt{ErrorBudget} & Deterministic three-term triangle inequality. \\
\texttt{Oracle} & Kernel axiom audit via \texttt{\#print axioms}. \\
\bottomrule
\end{tabular}
\caption{Current Lean 4 source proof boundary.}
\label{tab:lean}
\end{table}

The headline finite theorem has the schematic form:
\begin{lstlisting}[style=lean,language={}]
theorem hausdorffDiff_atomicMoment
    (weight point : iota -> Real) (k n : Nat) :
  hausdorffDiff k (atomicMoment weight point) n =
    sum i, weight i * point i ^ n * (1 - point i) ^ k
\end{lstlisting}
The repository CI is configured to run a full build, the oracle file, a placeholder scan, exact rational
certificate tests, and the paper build. No theorem in the current Lean source assumes RH.

\paragraph{Formalization frontier.}
The next high-value steps are the exact $\xi/\Xi$ convention bridge, the slit-plane map,
local factorization of analytic zeros, the Hausdorff theorem, and the prime-tail estimate.
The unbounded operator layer should be introduced only after this abstract chain is closed.

\section{Novelty, falsifiability, and publication gates}

The following distinctions are essential.

\paragraph{Classical substance.}
Stieltjes transforms, Hausdorff moments, normal families, and logarithmic-derivative pole
arguments are classical. Moment criteria related to RH already appear in the literature,
including \cite{Zhang2023}.

\paragraph{Potentially publishable package.}
A complete Lean formalization of \cref{thm:slit-extension,thm:one-point,prop:prime-tail},
combined with finite certified tests and a careful literature comparison, would be a
defensible formalization paper even without a new spectral convergence estimate.

\paragraph{Substantive analytic advance.}
A new unconditional bound for the Galerkin, ground-state alignment, model, or trace-resolvent
error of a concrete prime-built family would materially strengthen the project.

\paragraph{RH claim.}
Only a full, non-circular convergence theorem satisfying all domain, multiplicity, trace,
and normalization obligations would support a claim of proving RH. Finite numerical agreement,
no matter how accurate, is insufficient.

\section{Conclusion}

The squared-resolvent viewpoint turns a global zero-location problem into a sequence of
separable tasks. The key structural equivalence is
\[
 \boxed{
 \RH
 \iff
 \left\{
 x_0^n\frac{(-1)^n}{n!}\XiR^{(n)}(x_0)
 \right\}_{n\ge0}
 \text{ is a Hausdorff moment sequence}
 }
\]
for any fixed $x_0>1/4$. The forward direction compactifies the squared ordinates of the
zeros; the reverse direction uses a representing measure to build a holomorphic function on
the slit plane, where the logarithmic derivative cannot hide a non-real zero.

The resulting finite differences and semidefinite forms provide exact, falsifiable tests. The
prime expansion supplies a directly computable arithmetic target. Positive spectral approximants
need only be controlled at one point and converge on one real interval. What remains difficult is
not obscured: one must prove a uniform, non-circular comparison for a concrete prime-built
self-adjoint family. The repository is designed so that this open frontier is visible in the type and
dependency structure rather than disguised as an assumption named ``convergence.''

\appendix
\section{Finite atomic identity}

For completeness, let
\[
 b_n=\sum_{i=1}^mw_ip_i^n.
\]
Induction on $k$ proves
\[
 D^kb_n=\sum_{i=1}^mw_ip_i^n(1-p_i)^k.
\]
Indeed, the step from $k$ to $k+1$ is
\begin{align*}
 D^{k+1}b_n
 &=D^kb_n-D^kb_{n+1}\\
 &=\sum_iw_ip_i^n(1-p_i)^k
   -\sum_iw_ip_i^{n+1}(1-p_i)^k\\
 &=\sum_iw_ip_i^n(1-p_i)^{k+1}.
\end{align*}
This proof is mirrored directly in \texttt{PrimeResolvent/HausdorffFinite.lean}.

\section{Reproducibility commands}

From a fresh clone:
\begin{lstlisting}[style=lean,language={}]
lake exe cache get
lake build
lake env lean PrimeResolvent/Oracle.lean
python3 -m pip install -r requirements.txt
python3 -m pytest -q
python3 scripts/exact_atomic_certificate.py --output-dir data/certificates
python3 scripts/numerical_demo.py --output-dir data/demo
./scripts/build_paper.sh
\end{lstlisting}
The release manifest records SHA-256 hashes of every distributed source and generated artifact.
The numerical scripts are explicitly marked non-probative; exact toy certificates use rational
arithmetic.

\begin{thebibliography}{99}

\bibitem{Riemann1859}
B. Riemann,
\newblock Ueber die Anzahl der Primzahlen unter einer gegebenen Grösse,
\newblock \emph{Monatsberichte der Berliner Akademie} (1859), 671--680.

\bibitem{Titchmarsh1986}
E. C. Titchmarsh and D. R. Heath-Brown,
\newblock \emph{The Theory of the Riemann Zeta-Function}, 2nd ed.,
\newblock Oxford University Press, 1986.

\bibitem{Edwards1974}
H. M. Edwards,
\newblock \emph{Riemann's Zeta Function},
\newblock Academic Press, 1974.

\bibitem{Hausdorff1921}
F. Hausdorff,
\newblock Summationsmethoden und Momentfolgen. I,
\newblock \emph{Math. Z.} 9 (1921), 74--109.

\bibitem{ShohatTamarkin1943}
J. A. Shohat and J. D. Tamarkin,
\newblock \emph{The Problem of Moments},
\newblock American Mathematical Society, 1943.

\bibitem{Akhiezer1965}
N. I. Akhiezer,
\newblock \emph{The Classical Moment Problem},
\newblock Oliver and Boyd, 1965.

\bibitem{Widder1941}
D. V. Widder,
\newblock \emph{The Laplace Transform},
\newblock Princeton University Press, 1941.

\bibitem{Zhang2023}
R. Zhang,
\newblock An Application of Hausdorff Moment Problem,
\newblock arXiv:2303.09396, 2023.

\bibitem{ConnesConsaniMoscovici2025}
A. Connes, C. Consani, and H. Moscovici,
\newblock Zeta Spectral Triples,
\newblock arXiv:2511.22755, 2025.

\bibitem{ConnesVanSuijlekom2025}
A. Connes and W. D. van Suijlekom,
\newblock Quadratic Forms, Real Zeros and Echoes of the Spectral Action,
\newblock arXiv:2511.23257, 2025.

\bibitem{deMoura2015}
L. de Moura, S. Kong, J. Avigad, F. van Doorn, and J. von Raumer,
\newblock The Lean theorem prover (system description),
\newblock in \emph{CADE-25}, 2015, 378--388.

\bibitem{Mathlib2020}
The mathlib community,
\newblock The Lean mathematical library,
\newblock in \emph{CPP 2020}, 367--381.

\bibitem{ClayRH}
Clay Mathematics Institute,
\newblock Riemann Hypothesis, Millennium Prize Problems,
\newblock \url{https://www.claymath.org/millennium/riemann-hypothesis/}, accessed 22 June 2026.

\end{thebibliography}

\end{document}
